% Options for packages loaded elsewhere
\PassOptionsToPackage{unicode}{hyperref}
\PassOptionsToPackage{hyphens}{url}
\PassOptionsToPackage{dvipsnames,svgnames,x11names}{xcolor}
\documentclass[
  brazilian,
  11pt,
]{article}
\usepackage{xcolor}
\usepackage[margin=2.4cm]{geometry}
\usepackage{amsmath,amssymb}
\setcounter{secnumdepth}{-\maxdimen} % remove section numbering
\usepackage{iftex}
\ifPDFTeX
  \usepackage[T1]{fontenc}
  \usepackage[utf8]{inputenc}
  \usepackage{textcomp} % provide euro and other symbols
\else % if luatex or xetex
  \usepackage{unicode-math} % this also loads fontspec
  \defaultfontfeatures{Scale=MatchLowercase}
  \defaultfontfeatures[\rmfamily]{Ligatures=TeX,Scale=1}
\fi
\usepackage{lmodern}
\ifPDFTeX\else
  % xetex/luatex font selection
\fi
% Use upquote if available, for straight quotes in verbatim environments
\IfFileExists{upquote.sty}{\usepackage{upquote}}{}
\IfFileExists{microtype.sty}{% use microtype if available
  \usepackage[]{microtype}
  \UseMicrotypeSet[protrusion]{basicmath} % disable protrusion for tt fonts
}{}
\makeatletter
\@ifundefined{KOMAClassName}{% if non-KOMA class
  \IfFileExists{parskip.sty}{%
    \usepackage{parskip}
  }{% else
    \setlength{\parindent}{0pt}
    \setlength{\parskip}{6pt plus 2pt minus 1pt}}
}{% if KOMA class
  \KOMAoptions{parskip=half}}
\makeatother
\usepackage{color}
\usepackage{fancyvrb}
\newcommand{\VerbBar}{|}
\newcommand{\VERB}{\Verb[commandchars=\\\{\}]}
\DefineVerbatimEnvironment{Highlighting}{Verbatim}{commandchars=\\\{\}}
% Add ',fontsize=\small' for more characters per line
\usepackage{framed}
\definecolor{shadecolor}{RGB}{35,38,41}
\newenvironment{Shaded}{\begin{snugshade}}{\end{snugshade}}
\newcommand{\AlertTok}[1]{\textcolor[rgb]{0.58,0.85,0.30}{\textbf{\colorbox[rgb]{0.30,0.12,0.14}{#1}}}}
\newcommand{\AnnotationTok}[1]{\textcolor[rgb]{0.25,0.50,0.35}{#1}}
\newcommand{\AttributeTok}[1]{\textcolor[rgb]{0.16,0.50,0.73}{#1}}
\newcommand{\BaseNTok}[1]{\textcolor[rgb]{0.96,0.45,0.00}{#1}}
\newcommand{\BuiltInTok}[1]{\textcolor[rgb]{0.50,0.55,0.55}{#1}}
\newcommand{\CharTok}[1]{\textcolor[rgb]{0.24,0.68,0.91}{#1}}
\newcommand{\CommentTok}[1]{\textcolor[rgb]{0.48,0.49,0.49}{#1}}
\newcommand{\CommentVarTok}[1]{\textcolor[rgb]{0.50,0.55,0.55}{#1}}
\newcommand{\ConstantTok}[1]{\textcolor[rgb]{0.15,0.68,0.68}{\textbf{#1}}}
\newcommand{\ControlFlowTok}[1]{\textcolor[rgb]{0.99,0.74,0.29}{\textbf{#1}}}
\newcommand{\DataTypeTok}[1]{\textcolor[rgb]{0.16,0.50,0.73}{#1}}
\newcommand{\DecValTok}[1]{\textcolor[rgb]{0.96,0.45,0.00}{#1}}
\newcommand{\DocumentationTok}[1]{\textcolor[rgb]{0.64,0.20,0.25}{#1}}
\newcommand{\ErrorTok}[1]{\textcolor[rgb]{0.85,0.27,0.33}{\underline{#1}}}
\newcommand{\ExtensionTok}[1]{\textcolor[rgb]{0.00,0.60,1.00}{\textbf{#1}}}
\newcommand{\FloatTok}[1]{\textcolor[rgb]{0.96,0.45,0.00}{#1}}
\newcommand{\FunctionTok}[1]{\textcolor[rgb]{0.56,0.27,0.68}{#1}}
\newcommand{\ImportTok}[1]{\textcolor[rgb]{0.15,0.68,0.38}{#1}}
\newcommand{\InformationTok}[1]{\textcolor[rgb]{0.77,0.36,0.00}{#1}}
\newcommand{\KeywordTok}[1]{\textcolor[rgb]{0.81,0.81,0.76}{\textbf{#1}}}
\newcommand{\NormalTok}[1]{\textcolor[rgb]{0.81,0.81,0.76}{#1}}
\newcommand{\OperatorTok}[1]{\textcolor[rgb]{0.81,0.81,0.76}{#1}}
\newcommand{\OtherTok}[1]{\textcolor[rgb]{0.15,0.68,0.38}{#1}}
\newcommand{\PreprocessorTok}[1]{\textcolor[rgb]{0.15,0.68,0.38}{#1}}
\newcommand{\RegionMarkerTok}[1]{\textcolor[rgb]{0.16,0.50,0.73}{\colorbox[rgb]{0.08,0.19,0.26}{#1}}}
\newcommand{\SpecialCharTok}[1]{\textcolor[rgb]{0.24,0.68,0.91}{#1}}
\newcommand{\SpecialStringTok}[1]{\textcolor[rgb]{0.85,0.27,0.33}{#1}}
\newcommand{\StringTok}[1]{\textcolor[rgb]{0.96,0.31,0.31}{#1}}
\newcommand{\VariableTok}[1]{\textcolor[rgb]{0.15,0.68,0.68}{#1}}
\newcommand{\VerbatimStringTok}[1]{\textcolor[rgb]{0.85,0.27,0.33}{#1}}
\newcommand{\WarningTok}[1]{\textcolor[rgb]{0.85,0.27,0.33}{#1}}
\usepackage{longtable,booktabs,array}
\usepackage{caption}
\captionsetup[table]{skip=6pt}
\newcounter{none} % for unnumbered tables
\usepackage{calc} % for calculating minipage widths
% Correct order of tables after \paragraph or \subparagraph
\usepackage{etoolbox}
\makeatletter
\patchcmd\longtable{\par}{\if@noskipsec\mbox{}\fi\par}{}{}
\makeatother
% Allow footnotes in longtable head/foot
\IfFileExists{footnotehyper.sty}{\usepackage{footnotehyper}}{\usepackage{footnote}}
\makesavenoteenv{longtable}
\ifLuaTeX
\usepackage[bidi=basic,shorthands=off]{babel}
\else
\usepackage[bidi=default,shorthands=off]{babel}
\fi
\ifLuaTeX
  \usepackage{selnolig} % disable illegal ligatures
\fi
\setlength{\emergencystretch}{3em} % prevent overfull lines
\providecommand{\tightlist}{%
  \setlength{\itemsep}{0pt}\setlength{\parskip}{0pt}}
% ==========================================================================
%  Cabecalho LaTeX do curso "Trabalhando com Texto em Python I"
%  Injetado pelo pandoc (-H) ao gerar o .tex de cada unidade.
% ==========================================================================
\usepackage[T1]{fontenc}
\usepackage{lmodern}
\usepackage[scaled=0.92]{helvet}
\renewcommand{\familydefault}{\sfdefault}
\usepackage{microtype}
\usepackage{amssymb}
\usepackage{newunicodechar}
\usepackage{xcolor}
\usepackage[most]{tcolorbox}
\tcbuselibrary{skins,breakable}
\usepackage{titlesec}
\usepackage{enumitem}

% -------- paleta --------
\definecolor{accent}{HTML}{6C4CF1}
\definecolor{cyan}{HTML}{00B4D8}
\definecolor{subink}{HTML}{2B3242}
\definecolor{notec}{HTML}{3B82F6}
\definecolor{tipc}{HTML}{10B981}
\definecolor{warnc}{HTML}{F59E0B}

% -------- emojis / simbolos unicode que podem aparecer no texto --------
% simbolos com significado
\newunicodechar{✅}{{\color{tipc}\checkmark}}
\newunicodechar{✔}{{\color{tipc}\checkmark}}
\newunicodechar{☐}{$\square$}
\newunicodechar{κ}{\ensuremath{\kappa}}
\newunicodechar{≈}{\ensuremath{\approx}}
\newunicodechar{−}{\ensuremath{-}}
\newunicodechar{↓}{\ensuremath{\downarrow}}
\newunicodechar{→}{\ensuremath{\rightarrow}}
\newunicodechar{←}{\ensuremath{\leftarrow}}
\newunicodechar{↹}{\ensuremath{\hookrightarrow}}
\newunicodechar{°}{\textdegree}
\newunicodechar{·}{\textperiodcentered}
\newunicodechar{´}{'}
\newunicodechar{‑}{-}
% emojis decorativos -> removidos no PDF
\newunicodechar{💡}{}
\newunicodechar{🚀}{}
\newunicodechar{📝}{}
\newunicodechar{⚠}{}
\newunicodechar{🐍}{}
\newunicodechar{🎉}{}
\newunicodechar{🍕}{}
\newunicodechar{👍}{}
\newunicodechar{💋}{}
\newunicodechar{😚}{}
\newunicodechar{🚂}{}
\newunicodechar{❤}{}

% -------- caixas de callout (usadas pelo filtro callouts.lua) --------
\newtcolorbox{notebox}{enhanced,breakable,colback=notec!6,colframe=notec,
  boxrule=0pt,leftrule=4pt,arc=3pt,left=10pt,right=10pt,top=7pt,bottom=7pt,
  before skip=10pt,after skip=12pt}
\newtcolorbox{tipbox}{enhanced,breakable,colback=tipc!7,colframe=tipc,
  boxrule=0pt,leftrule=4pt,arc=3pt,left=10pt,right=10pt,top=7pt,bottom=7pt,
  before skip=10pt,after skip=12pt}
\newtcolorbox{warnbox}{enhanced,breakable,colback=warnc!8,colframe=warnc,
  boxrule=0pt,leftrule=4pt,arc=3pt,left=10pt,right=10pt,top=7pt,bottom=7pt,
  before skip=10pt,after skip=12pt}
\newtcolorbox{objbox}{enhanced,breakable,colback=accent!5,colframe=accent,
  boxrule=0pt,leftrule=4pt,arc=3pt,left=10pt,right=10pt,top=7pt,bottom=8pt,
  before skip=10pt,after skip=14pt}

% -------- titulos coloridos --------
\titleformat{\section}{\sffamily\Large\bfseries\color{accent}}{\thesection}{0.6em}{}
\titleformat{\subsection}{\sffamily\large\bfseries\color{subink}}{\thesubsection}{0.5em}{}
\titleformat{\subsubsection}{\sffamily\normalsize\bfseries\color{cyan}}{\thesubsubsection}{0.5em}{}
\titlespacing*{\section}{0pt}{2.2ex plus 1ex minus .2ex}{1.1ex}

% codigo inline colorido
\definecolor{inlink}{HTML}{5A34D6}
\let\oldtexttt\texttt
\renewcommand{\texttt}[1]{{\color{inlink}\oldtexttt{#1}}}
\usepackage{bookmark}
\IfFileExists{xurl.sty}{\usepackage{xurl}}{} % add URL line breaks if available
\urlstyle{same}
% fallback for those not using the hyperref driver hyperxmp:
\makeatletter
\@ifundefined{xmpquote}{\newcommand{\xmpquote}[1]{#1}}{}
\makeatother
\hypersetup{
  pdftitle={Avaliando modelos de linguagem},
  pdfauthor={CiberExt 26-29 · FEELT38103 · Universidade Federal de Uberlândia},
  pdflang={pt-BR},
  colorlinks=true,
  linkcolor={accent},
  filecolor={Maroon},
  citecolor={Blue},
  urlcolor={cyan},
  pdfcreator={LaTeX via pandoc}}

\title{Avaliando modelos de linguagem}
\usepackage{etoolbox}
\makeatletter
\providecommand{\subtitle}[1]{% add subtitle to \maketitle
  \apptocmd{\@title}{\par {\large #1 \par}}{}{}
}
\makeatother
\subtitle{Trabalhando com Texto em Python I · Unidade 4}
\author{CiberExt 26-29 · FEELT38103 · Universidade Federal de
Uberlândia}
\date{2026}

\begin{document}
\maketitle

{
\setcounter{tocdepth}{2}
\tableofcontents
}
\subsubsection{Objetivos de
aprendizagem}\label{objetivos-de-aprendizagem}

Ao final desta unidade, você deverá:

\begin{itemize}
\tightlist
\item
  entender o que é um \textbf{padrão‑ouro} (\emph{gold standard});
\item
  saber \textbf{avaliar a concordância} entre anotadores humanos;
\item
  compreender \textbf{métricas simples} para avaliar o desempenho do
  processamento de linguagem natural.
\end{itemize}

\begin{center}\rule{0.5\linewidth}{0.5pt}\end{center}

\subsection{\texorpdfstring{1. O que é um padrão‑ouro (\emph{gold
standard})?}{1. O que é um padrão‑ouro (gold standard)?}}\label{o-que-uxe9-um-padruxe3oouro-gold-standard}

Um \textbf{padrão‑ouro} (\emph{gold standard}) --- também chamado de
\emph{ground truth} (``verdade fundamental'') --- refere‑se a
\textbf{dados verificados por humanos} que podem ser usados como
referência (\emph{benchmark}) para avaliar o desempenho de algoritmos.
Em processamento de linguagem natural, os padrões‑ouro medem quão bem os
\textbf{humanos} desempenham uma tarefa.

O objetivo do PLN é permitir que computadores \textbf{alcancem ou
superem} o desempenho humano em uma tarefa pré‑definida. Medir se os
algoritmos conseguem isso exige uma referência --- e é o padrão‑ouro que
a fornece. Em termos simples, ele oferece um \textbf{ponto de
referência}.

É importante entender, porém, que \textbf{padrões‑ouro são abstrações do
uso da língua}. Considere a tarefa de classificar palavras em classes
gramaticais: as classes não são dadas pela natureza --- elas representam
uma abstração que impõe estrutura à língua. A língua, entretanto, é
naturalmente \textbf{ambígua e subjetiva}, e as abstrações usadas podem
ser incompletas: não temos certeza de que todos os usuários
classificariam as palavras da mesma forma.

Por isso precisamos \textbf{medir a confiabilidade} de qualquer
padrão‑ouro, ou seja, até que ponto os humanos \textbf{concordam} na
tarefa.

\begin{center}\rule{0.5\linewidth}{0.5pt}\end{center}

\subsection{2. Medindo a confiabilidade
manualmente}\label{medindo-a-confiabilidade-manualmente}

Esta seção mostra como a \textbf{confiabilidade} --- geralmente
entendida como a \textbf{concordância entre múltiplos anotadores} ---
pode ser medida manualmente.

\subsubsection{Passo 1 --- Anotar os
dados}\label{passo-1-anotar-os-dados}

A \textbf{análise de sentimento} é uma tarefa que consiste em determinar
o sentimento de um texto. Treinar um modelo de análise de sentimento
exige coletar dados de treino, isto é, exemplos de textos associados a
diferentes sentimentos.

Classifique os \emph{tweets} a seguir em três categorias ---
\textbf{positivo}, \textbf{neutro} ou \textbf{negativo} --- de acordo
com o sentimento. Anote sua decisão (uma por linha), mas \textbf{não}
discuta nem mostre ao colega ao lado:

\begin{enumerate}
\def\labelenumi{\arabic{enumi}.}
\tightlist
\item
  Updated: HSL GTFS (Helsinki, Finland) https://t.co/fWEpzmNQLz
\item
  current weather in Helsinki: broken clouds, -8°C 100\% humidity, wind
  4kmh, pressure 1061mb
\item
  CNN: ``WallStreetBets Redditors go ballistic over GameStop's sinking
  share price''
\item
  Baana bicycle counter. Today: 3 Same time last week: 1058 Trend: ↓99\%
  This year: 819 518 Last year: 802 079 \#Helsinki \#cycling
\item
  Elon Musk is now tweeting about \#bitcoin
\item
  A perfect Sunday walk in the woods just a few steps from home.
\item
  Went to Domino's today👍 It was so amazing and I think I got damn good
  dessert as well\ldots{}
\item
  Choo Choo 🚂 There's our train! 🎉 \#holidayahead
\item
  Happy women's day ❤💋 kisses to all you beautiful ladies. 😚
  \#awesometobeawoman
\item
  Good morning \#Helsinki! Sun will rise in 30 minutes (local time
  07:28)
\end{enumerate}

Anote suas classificações (uma por linha, de 1 a 10). O colega ao lado
faz o mesmo, de forma independente --- assim teremos \textbf{dois
anotadores} para comparar.

\subsubsection{Passo 2 --- Calcular a concordância
percentual}\label{passo-2-calcular-a-concorduxe2ncia-percentual}

Ao criar conjuntos de dados para treinar modelos, geralmente queremos
que os dados sejam \textbf{confiáveis}, ou seja, que haja concordância
sobre o que descrevemos (aqui, o sentimento dos \emph{tweets}).

Uma forma de medir isso é a \textbf{concordância percentual} simples:
quantas vezes, das 10, você e o colega concordaram. Basta dividir o
número de concordâncias pelo número de itens (10):

\begin{Shaded}
\begin{Highlighting}[]
\CommentTok{\# Substitua o numero abaixo pelo numero de itens em que voces concordaram}
\NormalTok{agreement }\OperatorTok{=} \DecValTok{0}

\CommentTok{\# Divide a contagem pelo numero de tweets}
\NormalTok{agreement }\OperatorTok{=}\NormalTok{ agreement }\OperatorTok{/} \DecValTok{10}

\CommentTok{\# Imprime a variavel}
\NormalTok{agreement}
\end{Highlighting}
\end{Shaded}

\begin{Shaded}
\begin{Highlighting}[]
\NormalTok{0.0}
\end{Highlighting}
\end{Shaded}

\subsubsection{Passo 3 --- Calcular as probabilidades de cada
categoria}\label{passo-3-calcular-as-probabilidades-de-cada-categoria}

A concordância percentual é, na verdade, uma medida \textbf{muito ruim}:
qualquer um dos dois pode ter \textbf{acertado por sorte} --- ou ter
achado a tarefa tediosa e classificado tudo aleatoriamente. Se isso
aconteceu, a concordância percentual \textbf{não tem como detectar},
pois ela não distingue acerto real de acaso!

Felizmente, dá para estimar a possibilidade de \textbf{concordância por
acaso}. O primeiro passo é contar quantas vezes \textbf{você} usou cada
categoria e converter essas contagens em \textbf{probabilidades},
dividindo pelo total de \emph{tweets}:

\begin{Shaded}
\begin{Highlighting}[]
\CommentTok{\# Conte quantos itens *voce* colocou em cada categoria}
\NormalTok{positive }\OperatorTok{=} \DecValTok{0}
\NormalTok{neutral }\OperatorTok{=} \DecValTok{0}
\NormalTok{negative }\OperatorTok{=} \DecValTok{0}

\CommentTok{\# Converte as contagens em probabilidades}
\NormalTok{positive }\OperatorTok{=}\NormalTok{ positive }\OperatorTok{/} \DecValTok{10}
\NormalTok{neutral }\OperatorTok{=}\NormalTok{ neutral }\OperatorTok{/} \DecValTok{10}
\NormalTok{negative }\OperatorTok{=}\NormalTok{ negative }\OperatorTok{/} \DecValTok{10}

\CommentTok{\# Chama cada variavel para examinar a saida}
\NormalTok{positive, neutral, negative}
\end{Highlighting}
\end{Shaded}

\begin{Shaded}
\begin{Highlighting}[]
\NormalTok{(0.0, 0.0, 0.0)}
\end{Highlighting}
\end{Shaded}

Essas probabilidades representam a chance de \textbf{você} escolher
aquela categoria. Agora pergunte ao colega as probabilidades dele e
registre‑as:

\begin{Shaded}
\begin{Highlighting}[]
\NormalTok{nb\_positive }\OperatorTok{=} \DecValTok{0}
\NormalTok{nb\_neutral }\OperatorTok{=} \DecValTok{0}
\NormalTok{nb\_negative }\OperatorTok{=} \DecValTok{0}
\end{Highlighting}
\end{Shaded}

Conhecendo as probabilidades de cada classe para os \textbf{dois}
anotadores, podemos calcular a probabilidade de ambos escolherem a mesma
categoria \textbf{por acaso}. Para cada categoria, basta multiplicar a
sua probabilidade pela do colega:

\begin{Shaded}
\begin{Highlighting}[]
\NormalTok{both\_positive }\OperatorTok{=}\NormalTok{ positive }\OperatorTok{*}\NormalTok{ nb\_positive}
\NormalTok{both\_neutral }\OperatorTok{=}\NormalTok{ neutral }\OperatorTok{*}\NormalTok{ nb\_neutral}
\NormalTok{both\_negative }\OperatorTok{=}\NormalTok{ negative }\OperatorTok{*}\NormalTok{ nb\_negative}
\end{Highlighting}
\end{Shaded}

\begin{notebox}

\textbf{Nota:} se um anotador não colocou nenhum \emph{tweet} em uma
categoria (por exemplo, \texttt{negative}) e o outro colocou, isso
\textbf{anula} a chance de concordância por acaso naquela categoria ---
multiplicar por zero resulta em zero.

\end{notebox}

\subsubsection{Passo 4 --- Estimar a concordância esperada (por
acaso)}\label{passo-4-estimar-a-concorduxe2ncia-esperada-por-acaso}

Agora podemos calcular quão provável é concordar \textbf{por acaso}.
Isso é a \textbf{concordância esperada} (\emph{expected agreement}),
obtida somando as probabilidades combinadas de cada categoria:

\begin{Shaded}
\begin{Highlighting}[]
\NormalTok{expected\_agreement }\OperatorTok{=}\NormalTok{ both\_positive }\OperatorTok{+}\NormalTok{ both\_neutral }\OperatorTok{+}\NormalTok{ both\_negative}

\NormalTok{expected\_agreement}
\end{Highlighting}
\end{Shaded}

\begin{Shaded}
\begin{Highlighting}[]
\NormalTok{0.0}
\end{Highlighting}
\end{Shaded}

Conhecendo tanto a \textbf{concordância observada} (\texttt{agreement})
quanto a \textbf{esperada por acaso} (\texttt{expected\_agreement}),
podemos usar uma medida mais confiável: o \textbf{kappa de Cohen} (κ),
que estima a concordância a partir de ambas. A fórmula é:

\[\kappa = \frac{P_{observado} - P_{esperado}}{1 - P_{esperado}}\]

Como essa informação já está nas variáveis \texttt{agreement} e
\texttt{expected\_agreement}, calculamos κ facilmente. Note que
envolvemos as subtrações em \textbf{parênteses} para executá‑las antes
da divisão:

\begin{Shaded}
\begin{Highlighting}[]
\NormalTok{kappa }\OperatorTok{=}\NormalTok{ (agreement }\OperatorTok{{-}}\NormalTok{ expected\_agreement) }\OperatorTok{/}\NormalTok{ (}\DecValTok{1} \OperatorTok{{-}}\NormalTok{ expected\_agreement)}

\NormalTok{kappa}
\end{Highlighting}
\end{Shaded}

\begin{Shaded}
\begin{Highlighting}[]
\NormalTok{0.0}
\end{Highlighting}
\end{Shaded}

\begin{center}\rule{0.5\linewidth}{0.5pt}\end{center}

\subsection{3. O kappa de Cohen como medida de
concordância}\label{o-kappa-de-cohen-como-medida-de-concorduxe2ncia}

O valor teórico do κ de Cohen vai de \textbf{−1} (discordância perfeita)
a \textbf{+1} (concordância perfeita), com \textbf{0} indicando
concordância totalmente aleatória. O κ costuma ser interpretado como uma
medida da \textbf{força} da concordância.

Landis e Koch (1977) propuseram os \emph{benchmarks} abaixo --- que
devem ser levados \textbf{com cautela}, pois as divisões são
arbitrárias:

{\def\LTcaptype{none} % do not increment counter
\begin{longtable}[]{@{}ll@{}}
\toprule\noalign{}
Kappa de Cohen (κ) & Força da concordância \\
\midrule\noalign{}
\endhead
\bottomrule\noalign{}
\endlastfoot
\textless{} 0,00 & Pobre \\
0,00 -- 0,20 & Ligeira \\
0,21 -- 0,40 & Razoável \\
0,41 -- 0,60 & Moderada \\
0,61 -- 0,80 & Substancial \\
0,81 -- 1,00 & Quase perfeita \\
\end{longtable}
}

O κ de Cohen serve para medir a concordância entre \textbf{dois}
anotadores, e as categorias disponíveis devem ser fixadas de antemão.
Para mais de dois anotadores, usa‑se uma medida como o \textbf{κ de
Fleiss}.

Na prática, o κ de Cohen (e muitas outras medidas) já está implementado
em bibliotecas Python --- raramente é preciso calcular à mão. A
biblioteca \textbf{scikit‑learn} (\texttt{sklearn}), por exemplo, inclui
a função \texttt{cohen\_kappa\_score()}, que recebe \textbf{duas listas}
e calcula o κ entre elas:

\begin{Shaded}
\begin{Highlighting}[]
\CommentTok{\# Importa a funcao cohen\_kappa\_score do modulo \textquotesingle{}metrics\textquotesingle{} da scikit{-}learn}
\ImportTok{from}\NormalTok{ sklearn.metrics }\ImportTok{import}\NormalTok{ cohen\_kappa\_score}

\CommentTok{\# Define duas listas, \textquotesingle{}a1\textquotesingle{} e \textquotesingle{}a2\textquotesingle{}, com anotacoes de classes gramaticais}
\NormalTok{a1 }\OperatorTok{=}\NormalTok{ [}\StringTok{\textquotesingle{}ADJ\textquotesingle{}}\NormalTok{, }\StringTok{\textquotesingle{}AUX\textquotesingle{}}\NormalTok{, }\StringTok{\textquotesingle{}NOUN\textquotesingle{}}\NormalTok{, }\StringTok{\textquotesingle{}VERB\textquotesingle{}}\NormalTok{, }\StringTok{\textquotesingle{}VERB\textquotesingle{}}\NormalTok{]}
\NormalTok{a2 }\OperatorTok{=}\NormalTok{ [}\StringTok{\textquotesingle{}ADJ\textquotesingle{}}\NormalTok{, }\StringTok{\textquotesingle{}VERB\textquotesingle{}}\NormalTok{, }\StringTok{\textquotesingle{}NOUN\textquotesingle{}}\NormalTok{, }\StringTok{\textquotesingle{}NOUN\textquotesingle{}}\NormalTok{, }\StringTok{\textquotesingle{}VERB\textquotesingle{}}\NormalTok{]}

\CommentTok{\# Usa cohen\_kappa\_score() para calcular a concordancia entre as listas}
\NormalTok{cohen\_kappa\_score(a1, a2)}
\end{Highlighting}
\end{Shaded}

\begin{Shaded}
\begin{Highlighting}[]
\NormalTok{0.44444444444444453}
\end{Highlighting}
\end{Shaded}

Segundo o \emph{benchmark} de Landis e Koch, esse valor indicaria
concordância \textbf{moderada}.

\begin{tipbox}

\textbf{Dica:} raramente é preciso anotar o conjunto inteiro para medir
a concordância --- uma \textbf{amostra aleatória} costuma bastar. Se o κ
sugere que os anotadores concordam, assumimos que as anotações são
confiáveis (não aleatórias). Ainda assim, toda medida de concordância
entre anotadores depende de suas próprias suposições sobre o que é
``concordar'' --- nenhuma representa a verdade absoluta.

\end{tipbox}

\begin{center}\rule{0.5\linewidth}{0.5pt}\end{center}

\subsection{4. Avaliando o desempenho de modelos de
linguagem}\label{avaliando-o-desempenho-de-modelos-de-linguagem}

Com um padrão‑ouro suficientemente confiável, podemos usá‑lo para medir
o desempenho de modelos de linguagem. Suponha um padrão‑ouro com 10
tokens anotados por classe gramatical, na lista \texttt{gold\_standard},
e as previsões de um modelo na lista \texttt{predictions}:

\begin{Shaded}
\begin{Highlighting}[]
\CommentTok{\# Define a lista \textquotesingle{}gold\_standard\textquotesingle{} (padrao{-}ouro anotado por humanos)}
\NormalTok{gold\_standard }\OperatorTok{=}\NormalTok{ [}\StringTok{\textquotesingle{}ADJ\textquotesingle{}}\NormalTok{, }\StringTok{\textquotesingle{}ADJ\textquotesingle{}}\NormalTok{, }\StringTok{\textquotesingle{}AUX\textquotesingle{}}\NormalTok{, }\StringTok{\textquotesingle{}VERB\textquotesingle{}}\NormalTok{, }\StringTok{\textquotesingle{}AUX\textquotesingle{}}\NormalTok{, }\StringTok{\textquotesingle{}NOUN\textquotesingle{}}\NormalTok{, }\StringTok{\textquotesingle{}NOUN\textquotesingle{}}\NormalTok{, }\StringTok{\textquotesingle{}ADJ\textquotesingle{}}\NormalTok{, }\StringTok{\textquotesingle{}DET\textquotesingle{}}\NormalTok{, }\StringTok{\textquotesingle{}PRON\textquotesingle{}}\NormalTok{]}

\CommentTok{\# Define a lista \textquotesingle{}predictions\textquotesingle{} (previsoes do modelo de linguagem)}
\NormalTok{predictions }\OperatorTok{=}\NormalTok{ [}\StringTok{\textquotesingle{}NOUN\textquotesingle{}}\NormalTok{, }\StringTok{\textquotesingle{}ADJ\textquotesingle{}}\NormalTok{, }\StringTok{\textquotesingle{}AUX\textquotesingle{}}\NormalTok{, }\StringTok{\textquotesingle{}VERB\textquotesingle{}}\NormalTok{, }\StringTok{\textquotesingle{}AUX\textquotesingle{}}\NormalTok{, }\StringTok{\textquotesingle{}NOUN\textquotesingle{}}\NormalTok{, }\StringTok{\textquotesingle{}VERB\textquotesingle{}}\NormalTok{, }\StringTok{\textquotesingle{}ADJ\textquotesingle{}}\NormalTok{, }\StringTok{\textquotesingle{}DET\textquotesingle{}}\NormalTok{, }\StringTok{\textquotesingle{}PROPN\textquotesingle{}}\NormalTok{]}
\end{Highlighting}
\end{Shaded}

\subsubsection{4.1 Acurácia}\label{acuruxe1cia}

Vamos importar o módulo \texttt{metrics} inteiro da scikit‑learn. A
função \texttt{accuracy\_score()} calcula a \textbf{acurácia}, que é
exatamente a mesma coisa que a concordância observada calculada
manualmente:

\begin{Shaded}
\begin{Highlighting}[]
\CommentTok{\# Importa o modulo \textquotesingle{}metrics\textquotesingle{} da biblioteca scikit{-}learn (sklearn)}
\ImportTok{from}\NormalTok{ sklearn }\ImportTok{import}\NormalTok{ metrics}

\CommentTok{\# Usa a funcao accuracy\_score() do modulo \textquotesingle{}metrics\textquotesingle{}}
\NormalTok{metrics.accuracy\_score(gold\_standard, predictions)}
\end{Highlighting}
\end{Shaded}

\begin{Shaded}
\begin{Highlighting}[]
\NormalTok{0.7}
\end{Highlighting}
\end{Shaded}

A acurácia sofre do mesmo problema da concordância observada --- pode
ser resultado de \textbf{palpites de sorte}. Como é um exemplo pequeno,
dá para verificar que \textbf{7 de 10} classes gramaticais coincidem,
resultando em acurácia de 0,7 (70\%).

\subsubsection{4.2 Matriz de confusão}\label{matriz-de-confusuxe3o}

Para avaliar melhor o modelo, organizamos os resultados em uma
\textbf{matriz de confusão}. Precisamos de todas as classes gramaticais
que aparecem em \texttt{gold\_standard} e \texttt{predictions}.
Coletamos as categorias únicas com a função \texttt{set()} --- um
\textbf{conjunto} (\emph{set}) é uma estrutura do Python com itens
únicos, útil aqui para remover duplicatas:

\begin{Shaded}
\begin{Highlighting}[]
\CommentTok{\# Coleta as classes gramaticais unicas em um conjunto, combinando as duas listas}
\NormalTok{pos\_tags }\OperatorTok{=} \BuiltInTok{set}\NormalTok{(gold\_standard }\OperatorTok{+}\NormalTok{ predictions)}

\CommentTok{\# Ordena o conjunto alfabeticamente e converte o resultado em lista}
\NormalTok{pos\_tags }\OperatorTok{=} \BuiltInTok{list}\NormalTok{(}\BuiltInTok{sorted}\NormalTok{(pos\_tags))}

\CommentTok{\# Imprime a lista resultante}
\NormalTok{pos\_tags}
\end{Highlighting}
\end{Shaded}

\begin{Shaded}
\begin{Highlighting}[]
\NormalTok{[\textquotesingle{}ADJ\textquotesingle{}, \textquotesingle{}AUX\textquotesingle{}, \textquotesingle{}DET\textquotesingle{}, \textquotesingle{}NOUN\textquotesingle{}, \textquotesingle{}PRON\textquotesingle{}, \textquotesingle{}PROPN\textquotesingle{}, \textquotesingle{}VERB\textquotesingle{}]}
\end{Highlighting}
\end{Shaded}

Usamos essas categorias para montar uma tabela em que as \textbf{linhas}
representam o padrão‑ouro e as \textbf{colunas} as previsões do modelo.
Percorremos cada par de itens (padrão‑ouro, previsão) e somamos
\texttt{+1} na célula correspondente. Por exemplo, o primeiro item de
\texttt{gold\_standard} é \texttt{ADJ} e o primeiro de
\texttt{predictions} é \texttt{NOUN}:

\begin{Shaded}
\begin{Highlighting}[]
\CommentTok{\# Imprime o primeiro item de cada lista}
\NormalTok{gold\_standard[}\DecValTok{0}\NormalTok{], predictions[}\DecValTok{0}\NormalTok{]}
\end{Highlighting}
\end{Shaded}

\begin{Shaded}
\begin{Highlighting}[]
\NormalTok{(\textquotesingle{}ADJ\textquotesingle{}, \textquotesingle{}NOUN\textquotesingle{})}
\end{Highlighting}
\end{Shaded}

Encontramos a linha \texttt{ADJ} e a coluna \texttt{NOUN} e somamos 1
àquela célula. A tabela completa (a \textbf{matriz de confusão}) fica
assim:

{\def\LTcaptype{none} % do not increment counter
\begin{longtable}[]{@{}lccccccc@{}}
\toprule\noalign{}
& ADJ & AUX & DET & NOUN & PRON & PROPN & VERB \\
\midrule\noalign{}
\endhead
\bottomrule\noalign{}
\endlastfoot
\textbf{ADJ} & 2 & 0 & 0 & 1 & 0 & 0 & 0 \\
\textbf{AUX} & 0 & 2 & 0 & 0 & 0 & 0 & 0 \\
\textbf{DET} & 0 & 0 & 1 & 0 & 0 & 0 & 0 \\
\textbf{NOUN} & 0 & 0 & 0 & 1 & 0 & 0 & 1 \\
\textbf{PRON} & 0 & 0 & 0 & 0 & 0 & 1 & 0 \\
\textbf{PROPN} & 0 & 0 & 0 & 0 & 0 & 0 & 0 \\
\textbf{VERB} & 0 & 0 & 0 & 0 & 0 & 0 & 1 \\
\end{longtable}
}

As previsões corretas formam uma \textbf{linha aproximadamente diagonal}
na tabela. A partir dela, derivamos duas métricas por classe:
\textbf{precisão} (\emph{precision}) e \textbf{revocação}
(\emph{recall}).

\begin{itemize}
\tightlist
\item
  \textbf{Precisão} é a proporção de previsões corretas \textbf{por
  classe} --- quantas das previsões daquela classe estavam certas. Por
  exemplo, a coluna \texttt{VERB} soma \texttt{2}, mas só \texttt{1}
  previsão está correta (na linha \texttt{VERB}); logo, a precisão de
  \texttt{VERB} é 1/2 = \textbf{0,5}. O mesmo vale para \texttt{NOUN}.
\item
  \textbf{Revocação} é a proporção de previsões corretas dentre
  \textbf{todos os exemplos reais} daquela classe --- quantas instâncias
  reais o modelo conseguiu ``encontrar''. Por exemplo, a linha
  \texttt{ADJ} soma \texttt{3} (há três adjetivos no padrão‑ouro), mas
  só \texttt{2} estão na coluna \texttt{ADJ}; logo, a revocação de
  \texttt{ADJ} é 2/3 ≈ \textbf{0,66}. Para \texttt{NOUN}, a revocação é
  1/2 = 0,5.
\end{itemize}

A scikit‑learn gera matrizes de confusão automaticamente com
\texttt{confusion\_matrix()}:

\begin{Shaded}
\begin{Highlighting}[]
\CommentTok{\# Calcula a matriz de confusao para as duas listas e imprime o resultado}
\BuiltInTok{print}\NormalTok{(metrics.confusion\_matrix(gold\_standard, predictions))}
\end{Highlighting}
\end{Shaded}

\begin{Shaded}
\begin{Highlighting}[]
\NormalTok{[[2 0 0 1 0 0 0]}
\NormalTok{ [0 2 0 0 0 0 0]}
\NormalTok{ [0 0 1 0 0 0 0]}
\NormalTok{ [0 0 0 1 0 0 1]}
\NormalTok{ [0 0 0 0 0 1 0]}
\NormalTok{ [0 0 0 0 0 0 0]}
\NormalTok{ [0 0 0 0 0 0 1]]}
\end{Highlighting}
\end{Shaded}

\subsubsection{4.3 Precisão e revocação com a
scikit‑learn}\label{precisuxe3o-e-revocauxe7uxe3o-com-a-scikitlearn}

A \textbf{precisão} é implementada na função
\texttt{precision\_score()}. Como temos \textbf{mais de duas} classes,
precisamos definir como os resultados de cada classe são combinados,
pelo argumento \texttt{average}. Com \texttt{average=None}, a função
calcula a precisão \textbf{de cada classe}. Também definimos
\texttt{zero\_division=0} para evitar erro quando uma classe de
\texttt{predictions} não existe em \texttt{gold\_standard} (nesses
casos, a precisão é 0):

\begin{Shaded}
\begin{Highlighting}[]
\CommentTok{\# Calcula a precisao entre as duas listas, para cada classe (classe gramatical)}
\NormalTok{precision }\OperatorTok{=}\NormalTok{ metrics.precision\_score(gold\_standard, predictions, average}\OperatorTok{=}\VariableTok{None}\NormalTok{, zero\_division}\OperatorTok{=}\DecValTok{0}\NormalTok{)}

\CommentTok{\# Chama a variavel para examinar o resultado}
\NormalTok{precision}
\end{Highlighting}
\end{Shaded}

\begin{Shaded}
\begin{Highlighting}[]
\NormalTok{array([1. , 1. , 1. , 0.5, 0. , 0. , 0.5])}
\end{Highlighting}
\end{Shaded}

A saída é um \textbf{array NumPy}. Para combinar os rótulos de
\texttt{pos\_tags} com as pontuações de \texttt{precision}, usamos a
função \texttt{zip()}, que une listas/arrays do mesmo tamanho; para
visualizar, convertemos com \texttt{dict()}:

\begin{Shaded}
\begin{Highlighting}[]
\CommentTok{\# Combina o conjunto \textquotesingle{}pos\_tags\textquotesingle{} com o array \textquotesingle{}precision\textquotesingle{} usando zip();}
\CommentTok{\# converte o resultado em um dicionario}
\BuiltInTok{dict}\NormalTok{(}\BuiltInTok{zip}\NormalTok{(pos\_tags, precision))}
\end{Highlighting}
\end{Shaded}

\begin{Shaded}
\begin{Highlighting}[]
\NormalTok{\{\textquotesingle{}ADJ\textquotesingle{}: 1.0, \textquotesingle{}AUX\textquotesingle{}: 1.0, \textquotesingle{}DET\textquotesingle{}: 1.0, \textquotesingle{}NOUN\textquotesingle{}: 0.5, \textquotesingle{}PRON\textquotesingle{}: 0.0, \textquotesingle{}PROPN\textquotesingle{}: 0.0, \textquotesingle{}VERB\textquotesingle{}: 0.5\}}
\end{Highlighting}
\end{Shaded}

Para uma \textbf{única} pontuação de precisão para todas as classes,
usamos \texttt{average=\textquotesingle{}macro\textquotesingle{}}, que
trata cada classe como igualmente importante, independentemente de
quantas instâncias ela tem:

\begin{Shaded}
\begin{Highlighting}[]
\CommentTok{\# Calcula a precisao entre as duas listas e tira a media (macro)}
\NormalTok{macro\_precision }\OperatorTok{=}\NormalTok{ metrics.precision\_score(gold\_standard, predictions, average}\OperatorTok{=}\StringTok{\textquotesingle{}macro\textquotesingle{}}\NormalTok{, zero\_division}\OperatorTok{=}\DecValTok{0}\NormalTok{)}

\CommentTok{\# Chama a variavel para examinar o resultado}
\NormalTok{macro\_precision}
\end{Highlighting}
\end{Shaded}

\begin{Shaded}
\begin{Highlighting}[]
\NormalTok{0.5714285714285714}
\end{Highlighting}
\end{Shaded}

A precisão \textbf{macro} é a soma das precisões dividida pelo número de
classes --- dá para verificar manualmente:

\begin{Shaded}
\begin{Highlighting}[]
\CommentTok{\# Calcula a media macro manualmente: soma as precisoes e divide}
\CommentTok{\# pelo numero de classes em \textquotesingle{}precision\textquotesingle{}}
\BuiltInTok{sum}\NormalTok{(precision) }\OperatorTok{/} \BuiltInTok{len}\NormalTok{(precision)}
\end{Highlighting}
\end{Shaded}

\begin{Shaded}
\begin{Highlighting}[]
\NormalTok{0.5714285714285714}
\end{Highlighting}
\end{Shaded}

A \textbf{revocação} é calculada do mesmo modo, com
\texttt{recall\_score()}:

\begin{Shaded}
\begin{Highlighting}[]
\CommentTok{\# Calcula a revocacao entre as duas listas, para cada classe}
\NormalTok{recall }\OperatorTok{=}\NormalTok{ metrics.recall\_score(gold\_standard, predictions, average}\OperatorTok{=}\VariableTok{None}\NormalTok{, zero\_division}\OperatorTok{=}\DecValTok{0}\NormalTok{)}

\CommentTok{\# Combina \textquotesingle{}pos\_tags\textquotesingle{} com o array \textquotesingle{}recall\textquotesingle{} e converte em dicionario}
\BuiltInTok{dict}\NormalTok{(}\BuiltInTok{zip}\NormalTok{(pos\_tags, recall))}
\end{Highlighting}
\end{Shaded}

\begin{Shaded}
\begin{Highlighting}[]
\NormalTok{\{\textquotesingle{}ADJ\textquotesingle{}: 0.6666666666666666, \textquotesingle{}AUX\textquotesingle{}: 1.0, \textquotesingle{}DET\textquotesingle{}: 1.0, \textquotesingle{}NOUN\textquotesingle{}: 0.5, \textquotesingle{}PRON\textquotesingle{}: 0.0, \textquotesingle{}PROPN\textquotesingle{}: 0.0, \textquotesingle{}VERB\textquotesingle{}: 1.0\}}
\end{Highlighting}
\end{Shaded}

\subsubsection{4.4 Relatório de
classificação}\label{relatuxf3rio-de-classificauxe7uxe3o}

A scikit‑learn oferece a função \texttt{classification\_report()}, que
dá precisão e revocação de cada classe, junto com o \textbf{F1‑score}
--- uma média equilibrada entre precisão e revocação (ambas contribuem
igualmente), variando de 0 a 1:

\begin{Shaded}
\begin{Highlighting}[]
\CommentTok{\# Imprime um relatorio de classificacao}
\BuiltInTok{print}\NormalTok{(metrics.classification\_report(gold\_standard, predictions, zero\_division}\OperatorTok{=}\DecValTok{0}\NormalTok{))}
\end{Highlighting}
\end{Shaded}

\begin{Shaded}
\begin{Highlighting}[]
\NormalTok{              precision    recall  f1{-}score   support}

\NormalTok{         ADJ       1.00      0.67      0.80         3}
\NormalTok{         AUX       1.00      1.00      1.00         2}
\NormalTok{         DET       1.00      1.00      1.00         1}
\NormalTok{        NOUN       0.50      0.50      0.50         2}
\NormalTok{        PRON       0.00      0.00      0.00         1}
\NormalTok{       PROPN       0.00      0.00      0.00         0}
\NormalTok{        VERB       0.50      1.00      0.67         1}

\NormalTok{    accuracy                           0.70        10}
\NormalTok{   macro avg       0.57      0.60      0.57        10}
\NormalTok{weighted avg       0.75      0.70      0.71        10}
\end{Highlighting}
\end{Shaded}

As pontuações da linha \texttt{macro\ avg} correspondem às que
calculamos acima. A linha \texttt{weighted\ avg} (média
\textbf{ponderada}) leva em conta o número de instâncias de cada classe;
a coluna \texttt{support} conta quantas instâncias foram observadas em
cada classe.

\begin{center}\rule{0.5\linewidth}{0.5pt}\end{center}

\subsection{Quiz}\label{quiz}

Marque a alternativa correta (a resposta certa está destacada com ✅).

\textbf{1. Um padrão‑ouro (\emph{gold standard}) é:}

\begin{enumerate}
\def\labelenumi{\arabic{enumi}.}
\tightlist
\item
  Dados verificados por humanos, usados como referência ✅
\item
  A saída bruta do modelo
\item
  Um dicionário do Python
\end{enumerate}

\textbf{2. Por que a concordância percentual é uma medida fraca?}

\begin{enumerate}
\def\labelenumi{\arabic{enumi}.}
\tightlist
\item
  Não distingue acerto real de acerto por acaso ✅
\item
  É difícil de calcular
\item
  Só funciona com duas classes
\end{enumerate}

\textbf{3. O kappa de Cohen mede a concordância entre quantos
anotadores?}

\begin{enumerate}
\def\labelenumi{\arabic{enumi}.}
\tightlist
\item
  Dois ✅
\item
  Três
\item
  Qualquer número
\end{enumerate}

\textbf{4. A precisão de uma classe responde: das vezes em que o modelo
\emph{previu} essa classe,}

\begin{enumerate}
\def\labelenumi{\arabic{enumi}.}
\tightlist
\item
  quantas estavam certas ✅
\item
  quantas existiam no total
\item
  quantas ele deixou de encontrar
\end{enumerate}

\textbf{5. A revocação de uma classe responde: dos exemplos que
\emph{realmente} são dessa classe,}

\begin{enumerate}
\def\labelenumi{\arabic{enumi}.}
\tightlist
\item
  quantos o modelo encontrou ✅
\item
  quantos ele previu a mais
\item
  quantos estavam errados
\end{enumerate}

\textbf{6. O F1‑score é:}

\begin{enumerate}
\def\labelenumi{\arabic{enumi}.}
\tightlist
\item
  A média equilibrada entre precisão e revocação ✅
\item
  A soma da precisão com a revocação
\item
  O maior valor entre os dois
\end{enumerate}

\subsection{Resumo da unidade}\label{resumo-da-unidade}

Nesta unidade, você aprendeu a:

\begin{enumerate}
\def\labelenumi{\arabic{enumi}.}
\tightlist
\item
  entender o que é um \textbf{padrão‑ouro} e por que ele é uma
  \textbf{abstração} do uso da língua;
\item
  medir a \textbf{confiabilidade} entre anotadores: concordância
  percentual, probabilidades por categoria e concordância esperada por
  acaso;
\item
  calcular e interpretar o \textbf{kappa de Cohen} (manualmente e com
  \texttt{cohen\_kappa\_score()}), usando os \emph{benchmarks} de Landis
  \& Koch;
\item
  avaliar o desempenho de modelos com \textbf{acurácia}, \textbf{matriz
  de confusão}, \textbf{precisão}, \textbf{revocação} e
  \textbf{F1‑score}, usando a \texttt{scikit‑learn}
  (\texttt{accuracy\_score}, \texttt{confusion\_matrix},
  \texttt{precision\_score}, \texttt{recall\_score},
  \texttt{classification\_report}).
\end{enumerate}

\subsubsection{Exercícios sugeridos}\label{exercuxedcios-sugeridos}

\begin{enumerate}
\def\labelenumi{\arabic{enumi}.}
\tightlist
\item
  Anote você mesmo os 10 \emph{tweets} e peça a um colega que faça o
  mesmo; calcule a concordância percentual e o κ de Cohen com
  \texttt{cohen\_kappa\_score()}.
\item
  Monte suas próprias listas \texttt{gold\_standard} e
  \texttt{predictions} (com pelo menos 8 itens) e gere a matriz de
  confusão com \texttt{confusion\_matrix()}.
\item
  A partir do relatório de classificação, explique por que a classe
  \texttt{PROPN} tem \texttt{support} igual a 0 e o que isso significa.
\item
  Calcule o F1‑score de \texttt{VERB} \textbf{manualmente} a partir da
  precisão e da revocação (fórmula: F1 = 2 · (P · R) / (P + R)) e
  compare com o relatório.
\end{enumerate}

\begin{center}\rule{0.5\linewidth}{0.5pt}\end{center}

\subsubsection{Referências}\label{referuxeancias}

\begin{itemize}
\tightlist
\item
  Rögnvaldsson, E. et al.~\emph{Applied Language Technology} (MOOC).
  Universidade de Helsinque.
  \url{https://applied-language-technology.mooc.fi/}
\item
  Landis, J. R.; Koch, G. G. (1977). The measurement of observer
  agreement for categorical data. \emph{Biometrics}.
\item
  Pedregosa, F. et al.~(2011). Scikit‑learn: Machine Learning in Python.
  \emph{JMLR}. \url{https://scikit-learn.org/}
\end{itemize}

\end{document}
